\documentclass[10pt]{article}

\usepackage[top=0.55in, bottom=0.65in, left=0.55in, right=0.55in]{geometry}
\usepackage{amsmath, amssymb}
\usepackage[T1]{fontenc}
\usepackage[utf8]{inputenc}
\usepackage{lmodern}
\usepackage{enumitem}
\usepackage{multicol}
\usepackage{hyperref}

\setlength{\parindent}{0pt}
\setlength{\parskip}{0.35em}
\setlength{\columnsep}{0.22in}
\setlist[itemize]{leftmargin=*, topsep=2pt, itemsep=1pt, parsep=0pt}
\setlist[enumerate]{leftmargin=*, topsep=2pt, itemsep=1pt, parsep=0pt}

\newcommand{\PPL}{\operatorname{PPL}}
\newcommand{\softmax}{\operatorname{softmax}}
\newcommand{\CE}{\operatorname{CE}}
\newenvironment{symdefs}{\begin{itemize}[leftmargin=1.1em, topsep=1pt, itemsep=0pt, parsep=0pt]}{\end{itemize}}
\newcommand{\sym}[2]{\item \textbf{#1}: #2}

\begin{document}
\small

\begin{center}
{\Large \textbf{6.8610 Quiz 1 Study Note}}\\
Spring 2026\\
Aggregated from Lectures 1--7 in \texttt{6.8610/LectureNotes}
\end{center}

\begin{multicols*}{2}
\raggedcolumns

\section*{Course Throughline}
\begin{itemize}
\item \textbf{Natural Language Processing (NLP)} = representation \(+\) modeling for human language.
\item Repeated theme: many NLP tasks can be reframed as language modeling.
\item Repeated warning: low loss or high fluency does \emph{not} imply true understanding, reliability, or usefulness.
\end{itemize}

\section*{Lecture 1: Language Models}
\textbf{Language model (LM):}
\[
p_\theta(x_{1:T}).
\]
\begin{symdefs}
\sym{\(x_{1:T}\)}{token sequence \((x_1,\dots,x_T)\)}
\sym{\(T\)}{sequence length}
\sym{\(\theta\)}{model parameters}
\end{symdefs}

\textbf{Autoregressive factorization:}
\[
p_\theta(x_{1:T})=\prod_{t=1}^{T} p_\theta(x_t \mid x_{1:t-1}).
\]

\textbf{Conditional language modeling:}
\[
p_\theta(y_{1:m}\mid x_{1:n})
=
\prod_{t=1}^{m} p_\theta(y_t \mid y_{1:t-1}, x_{1:n}).
\]
\begin{symdefs}
\sym{\(x_{1:n}\)}{input sequence}
\sym{\(y_{1:m}\)}{output sequence}
\sym{\(n\)}{input length}
\sym{\(m\)}{output length}
\end{symdefs}

\textbf{Maximum likelihood:}
\[
\hat{\theta}=\arg\max_\theta p_\theta(x_{1:T}).
\]

\textbf{Naive table size:}
\[
\sum_{t=1}^{T} V^t
\]
for vocabulary size \(V\); this is exponential in context length.

\textbf{Perplexity:}
\[
\PPL(x_{1:T})
=
\exp\left(
-\frac{1}{T}\log p_\theta(x_{1:T})
\right)
=
p_\theta(x_{1:T})^{-1/T}.
\]
\begin{symdefs}
\sym{\(\PPL\)}{per-token perplexity}
\sym{Lower \(\PPL\)}{better fit to the evaluation distribution}
\end{symdefs}

\begin{itemize}
\item Interpret perplexity as the model's average branching factor.
\item For a uniform distribution over \(V\) words, \(\PPL=V\).
\item \textbf{Perplexity sanity checks:} \(\PPL>0\) always, and a unigram maximum-likelihood model should outperform a uniform model on its own corpus.
\item \textbf{Shorter strings are not automatically better.} Autoregressive sentence probability also depends on whether the model decides to stop, so an end-of-sequence choice can make a longer complete sentence more probable than a shorter prefix.
\item \textbf{Conditional perplexity reaches \(1\) only for deterministic mappings.} If one input can map to multiple valid outputs, the conditional perplexity must stay above \(1\).
\item Language is hard because it is ambiguous, compositional, socially grounded, and constantly changing.
\item Representation and modeling are the two recurring pillars of the course.
\end{itemize}

\section*{Lecture 2: Log-Linear Models and Embeddings}
\textbf{Historical translation framing:}
\[
p(y \mid x) \propto p(x \mid y)\,p(y).
\]
\begin{symdefs}
\sym{\(x\)}{source sentence}
\sym{\(y\)}{target sentence}
\sym{\(p(y)\)}{target-language model}
\end{symdefs}

\textbf{Log-linear next-token model:}
\[
p_\theta(y \mid x)
=
\frac{\exp(\theta_y^\top \phi(x))}
{\sum_{y' \in V}\exp(\theta_{y'}^\top \phi(x))}.
\]
\begin{symdefs}
\sym{\(x\)}{context}
\sym{\(y\)}{next token}
\sym{\(\phi(x)\)}{fixed feature vector for context \(x\)}
\sym{\(\theta_y\)}{parameter vector for token \(y\)}
\sym{\(V\)}{vocabulary}
\end{symdefs}

\textbf{Single-example log-likelihood:}
\[
\log p_\theta(y \mid x)
=
\theta_y^\top \phi(x)-\log Z(x),
\qquad
Z(x)=\sum_{y' \in V}\exp(\theta_{y'}^\top \phi(x)).
\]
\begin{symdefs}
\sym{\(Z(x)\)}{partition function that normalizes scores into probabilities}
\end{symdefs}

\textbf{Gradient:}
\[
\nabla_{\theta_{y_0}}\log p_\theta(y \mid x)
=
\left(\mathbf{1}[y=y_0]-p_\theta(y_0 \mid x)\right)\phi(x).
\]
\begin{symdefs}
\sym{\(y_0\)}{candidate token}
\sym{\(\mathbf{1}[y=y_0]\)}{indicator that \(y_0\) is the gold token}
\end{symdefs}

\textbf{Term frequency-inverse document frequency (TF-IDF):}
\[
\operatorname{TF}(w,d)=\text{count of } w \text{ in } d,
\]
\[
\operatorname{IDF}(w)=\log \frac{N}{\operatorname{df}(w)},
\qquad
\operatorname{TFIDF}(w,d)=\operatorname{TF}(w,d)\operatorname{IDF}(w).
\]
\begin{symdefs}
\sym{\(d\)}{document}
\sym{\(w\)}{word type}
\sym{\(N\)}{number of documents}
\sym{\(\operatorname{df}(w)\)}{number of documents containing \(w\)}
\end{symdefs}

\textbf{Skip-gram:}
\[
p(y \mid x)
=
\frac{\exp(u_y^\top w_x)}
{\sum_{y' \in V}\exp(u_{y'}^\top w_x)}.
\]
\begin{symdefs}
\sym{\(w_x\)}{input embedding of center word \(x\)}
\sym{\(u_y\)}{output embedding of context word \(y\)}
\end{symdefs}
Predict context from center word.

\textbf{Continuous Bag-of-Words (CBOW):}
\[
W(x)=\frac{1}{|x|}\sum_{x_i \in x} w_{x_i},
\qquad
p(y \mid x)\propto \exp(u_y^\top W(x)).
\]
\begin{symdefs}
\sym{\(x\)}{context-word set or context window}
\sym{\(|x|\)}{number of context words}
\sym{\(W(x)\)}{averaged context embedding}
\end{symdefs}
Predict target word from neighboring words.

\textbf{Pointwise mutual information (PMI) view:}
\[
u_y^\top w_x \approx \log \frac{p(x,y)}{p(x)p(y)} - c.
\]
\begin{symdefs}
\sym{\(c\)}{approximately constant offset}
\end{symdefs}

\begin{itemize}
\item Bag-of-words and CBOW are order-insensitive.
\item Reversing the corpus should leave standard CBOW nearly unchanged because the local context multisets are preserved.
\item Skip-gram and CBOW learn useful \textbf{type-level} embeddings, not contextualized token representations.
\item Embedding quality can be checked intrinsically (similarity/analogy resources) or extrinsically (downstream-task performance).
\item CBOW and Skip-gram are \textbf{not} proper left-to-right language models.
\item Neural embedding methods are not guaranteed to beat count-based models on perplexity, because some of them are not autoregressive sequence models at all.
\item Embedding dimensions are usually not individually interpretable.
\item High input-output dot products such as \(u_y^\top w_x\) reflect predictive compatibility or co-occurrence association.
\item Similar contexts are more naturally reflected by similarity between embeddings of two words, not by a single cross dot product.
\item For two candidates \(a\) and \(b\) under the same softmax context,
\[
\frac{p(a \mid x)}{p(b \mid x)}=\exp(\operatorname{score}(a,x)-\operatorname{score}(b,x)).
\]
\item In word2vec-style models, the embedding layer is essentially a learned lookup table; richer contextual composition arrives in later sequence models.
\end{itemize}

\section*{Lecture 3: Feed-Forward Models, RNNs, Seq2Seq, Attention}
\textbf{Key contrast: masked vs autoregressive feed-forward models}
\begin{itemize}
\item \textbf{Masked objective:} predict a hidden token from surrounding context.
\item \textbf{Autoregressive objective:} predict the next token from left context only.
\item Masked objectives can learn representations; autoregressive objectives directly define normalized sequence probabilities.
\end{itemize}

\textbf{Basic recurrent neural network (RNN) update:}
\[
h_t = f(Wx_t + Vh_{t-1} + b_h),
\qquad
\hat{y}_t = \softmax(Uh_t + b_y).
\]
\begin{symdefs}
\sym{\(x_t\)}{input at timestep \(t\)}
\sym{\(h_t\)}{hidden state at timestep \(t\)}
\sym{\(\hat{y}_t\)}{predicted output distribution at timestep \(t\)}
\sym{\(\softmax(z)_j\)}{normalized exponential,
\[
\softmax(z)_j=\frac{e^{z_j}}{\sum_k e^{z_k}},
\]
where \(z\) is a vector of logits and \(j\) indexes one output coordinate}
\sym{\(W,V,U\)}{learned parameter matrices}
\sym{\(b_h,b_y\)}{bias terms}
\end{symdefs}

\textbf{Cross-entropy (CE) sequence loss:}
\[
\mathcal{L}=\sum_{t=1}^{T}\CE(y_t,\hat{y}_t),
\qquad
\CE(y_t,\hat{y}_t)= -\sum_{j \in V} y_{t,j}\log \hat{y}_{t,j}.
\]
\begin{symdefs}
\sym{\(y_t\)}{one-hot gold label at timestep \(t\)}
\sym{\(\hat{y}_{t,j}\)}{predicted probability of vocabulary item \(j\) at timestep \(t\)}
\sym{\(V\)}{vocabulary}
\end{symdefs}

\textbf{Backpropagation Through Time (BPTT), final-loss contribution:}
\[
\frac{\partial \mathcal{L}_T}{\partial h_k}
=
\frac{\partial \mathcal{L}_T}{\partial h_T}
\prod_{t=k+1}^{T}
\frac{\partial h_t}{\partial h_{t-1}}.
\]
\begin{symdefs}
\sym{\(\mathcal{L}_T\)}{loss contribution from the final timestep \(T\)}
\sym{Interpretation}{for full sequence loss, gradients to \(h_k\) sum contributions from later timesteps; repeated Jacobian products can make gradients vanish or explode}
\end{symdefs}

\textbf{Exploding-gradient fix: gradient clipping}
\[
g \leftarrow g \cdot \min\left(1,\frac{\tau}{\lVert g \rVert}\right).
\]
\begin{symdefs}
\sym{\(g\)}{gradient vector}
\sym{\(\lVert g \rVert\)}{gradient norm}
\sym{\(\tau\)}{clipping threshold}
\end{symdefs}

\textbf{Seq2Seq encoder-decoder:}
\[
h^{dec}_0 = h^{enc}_N,
\qquad
h^{dec}_t = f(W_y y_{t-1} + V h^{dec}_{t-1} + b).
\]
\begin{symdefs}
\sym{\(h_N^{enc}\)}{final encoder state}
\sym{\(h_t^{dec}\)}{decoder state at timestep \(t\)}
\sym{\(y_{t-1}\)}{previous target-token embedding or representation}
\end{symdefs}

\textbf{Attention score and weights:}
\[
e_{t,i}=\operatorname{score}(h_t^{dec}, h_i^{enc}),
\qquad
a_{t,i}=\frac{\exp(e_{t,i})}{\sum_{j=1}^{N}\exp(e_{t,j})}.
\]
\begin{symdefs}
\sym{\(t\)}{decoder timestep}
\sym{\(i\)}{source position}
\sym{\(e_{t,i}\)}{compatibility score between decoder step \(t\) and source position \(i\)}
\sym{\(a_{t,i}\)}{normalized attention weight}
\end{symdefs}

\textbf{Context vector:}
\[
c_t=\sum_{i=1}^{N} a_{t,i} h_i^{enc}.
\]
\begin{symdefs}
\sym{\(c_t\)}{context vector, i.e. weighted average of encoder states}
\end{symdefs}

\textbf{RNN can emulate a 2-token feed-forward LM:}
\[
h_t=[Wx_{t-1};\,Wx_{t-2}],
\]
so the hidden state stores the last two token embeddings and reproduces a fixed 2-token context window.
\begin{symdefs}
\sym{\(W\)}{embedding matrix for input tokens}
\sym{\(x_{t-1}, x_{t-2}\)}{one-hot vectors for the previous two input tokens}
\sym{\(h_t\)}{RNN hidden state storing the embeddings of the two most recent tokens}
\sym{\([a;b]\)}{concatenation of vector \(a\) above vector \(b\)}
\end{symdefs}

\textbf{Popular attention scoring functions:}
\begin{itemize}
\item Additive / concat:
\[
e_{t,i}=v_a^\top \tanh\!\left(W_a[h_t^{dec}; h_i^{enc}]\right).
\]
\item Bilinear / general:
\[
e_{t,i}=(h_t^{dec})^\top W_a h_i^{enc}.
\]
\item Dot product:
\[
e_{t,i}=(h_t^{dec})^\top h_i^{enc}.
\]
\end{itemize}

\begin{itemize}
\item RNN hidden states are \textbf{contextualized token embeddings}; word2vec embeddings are static type vectors.
\item Teacher forcing: during training, feed the gold previous token to the decoder.
\item Truncated BPTT truncates the \textbf{backward} pass, not the forward recurrent state.
\item RNN strengths: shared weights, variable-length processing, stateful memory.
\item RNN weaknesses: slow, hard to parallelize, vanishing/exploding gradients, weak practical long-range memory.
\item Long short-term memory (LSTM) networks are a gated repair to vanilla RNNs.
\item Attention breaks the single-vector encoder bottleneck and is the bridge to Transformers.
\end{itemize}

\section*{Lecture 4: Transformers}
\textbf{Queries / keys / values:}
\[
q_i=W_Q x_i,
\qquad
k_i=W_K x_i,
\qquad
v_i=W_V x_i.
\]
\begin{symdefs}
\sym{\(x_i\)}{representation of token \(i\)}
\sym{\(q_i\)}{query vector of token \(i\)}
\sym{\(k_i\)}{key vector of token \(i\)}
\sym{\(v_i\)}{value vector of token \(i\)}
\end{symdefs}

\textbf{Scaled dot-product attention:}
\[
a_{ij}
=
\softmax_j\!\left(\frac{q_i^\top k_j}{\sqrt{d_k}}\right),
\qquad
z_i=\sum_{j=1}^{n} a_{ij} v_j.
\]
\begin{symdefs}
\sym{\(a_{ij}\)}{attention weight from token \(i\) to token \(j\)}
\sym{\(d_k\)}{key dimension}
\sym{\(n\)}{sequence length}
\sym{\(z_i\)}{new representation for token \(i\)}
\end{symdefs}

\textbf{Matrix form:}
\[
\operatorname{Attention}(Q,K,V)
=
\softmax\!\left(\frac{QK^\top}{\sqrt{d_k}}\right)V.
\]
\begin{symdefs}
\sym{\(Q\)}{matrix of token queries}
\sym{\(K\)}{matrix of token keys}
\sym{\(V\)}{matrix of token values}
\end{symdefs}

\textbf{Multi-Head Attention (MHA):}
\[
\text{head}_h
=
\operatorname{Attention}(QW_Q^{(h)}, KW_K^{(h)}, VW_V^{(h)}),
\]
\[
\operatorname{MHA}(Q,K,V)
=
\operatorname{Concat}(\text{head}_1,\dots,\text{head}_H)W_O.
\]
\begin{symdefs}
\sym{\(H\)}{number of attention heads}
\sym{\(W_O\)}{output projection matrix}
\end{symdefs}

\textbf{Position-wise feed-forward neural network (FFNN):}
\[
\operatorname{FFNN}(z_i)=W_2\,\sigma(W_1 z_i + b_1)+b_2.
\]
\begin{symdefs}
\sym{\(\sigma\)}{position-wise nonlinearity}
\end{symdefs}

\textbf{Residual connection:}
\[
f_{\text{res}}(x)=f(x)+x.
\]

\textbf{Layer normalization:}
\[
\hat{z}_{t,j}=\frac{z_{t,j}-\mu_t}{\sigma_t+\epsilon},
\qquad
y_{t,j}=\gamma_j \hat{z}_{t,j}+\beta_j.
\]
\begin{symdefs}
\sym{\(\mu_t\)}{mean over features of token \(t\)}
\sym{\(\sigma_t\)}{standard deviation over features of token \(t\)}
\sym{\(\epsilon\)}{small stability constant}
\sym{\(\gamma,\beta\)}{learned scale and shift parameters}
\end{symdefs}

\textbf{Positional encodings:}
\[
x_i' = x_i + p_i.
\]
\begin{symdefs}
\sym{\(p_i\)}{position encoding for token \(i\)}
\end{symdefs}
Sinusoidal:
\[
\operatorname{PE}_{(\text{pos},2i)}
=
\sin\!\left(\frac{\text{pos}}{10000^{2i/d_{\text{model}}}}\right),
\]
\[
\operatorname{PE}_{(\text{pos},2i+1)}
=
\cos\!\left(\frac{\text{pos}}{10000^{2i/d_{\text{model}}}}\right).
\]

\textbf{Masked self-attention:}
\[
\operatorname{MaskedAttention}(Q,K,V)
=
\softmax\!\left(\frac{QK^\top + M}{\sqrt{d_k}}\right)V,
\]
with \(M_{ij}=0\) if \(j \le i\), else \(-\infty\).
\begin{symdefs}
\sym{\(M\)}{causal mask preventing token \(i\) from attending to future positions \(j>i\)}
\end{symdefs}

\textbf{Cross-attention:}
\[
\operatorname{CrossAttention}
=
\softmax\!\left(
\frac{Q^{dec}(K^{enc})^\top}{\sqrt{d_k}}
\right)V^{enc}.
\]
\begin{symdefs}
\sym{\(Q^{dec}\)}{queries from decoder states}
\sym{\(K^{enc}\)}{keys from encoder states}
\sym{\(V^{enc}\)}{values from encoder states}
\end{symdefs}

\textbf{BERT objectives:}
\[
\max_\theta \sum_{t \in \mathcal{M}} \log p_\theta(x_t \mid x_{\setminus \mathcal{M}}),
\qquad
p_\theta(\texttt{isNext} \mid s_1, s_2).
\]
\begin{symdefs}
\sym{\(\mathcal{M}\)}{set of masked token positions}
\sym{\(s_1,s_2\)}{sentence pair used in next-sentence prediction}
\end{symdefs}

\begin{itemize}
\item Three ways to attend:
\begin{itemize}
\item self-attention: one sequence attends to itself;
\item masked self-attention: same, but future positions blocked;
\item cross-attention: decoder attends to encoder states.
\end{itemize}
\item Self-attention lets each token access all other positions directly and parallelizes across positions better than recurrent models.
\item Without positional information, a Transformer encoder behaves like a bag-of-words system up to permutation symmetry.
\item Positional information is necessary because content-only attention does not otherwise encode order.
\item Transformer encoder = self-attention \(+\) residual/norm \(+\) FFNN.
\item Transformer decoder = masked self-attention; encoder-decoder models also add cross-attention for conditional generation.
\item Causal masking matters because it preserves the left-to-right training objective and prevents information leakage from future tokens.
\item \textbf{BERT} = encoder-only, pretrained with masked language modeling and originally next sentence prediction.
\item BERT masking rule: choose \(15\%\) of tokens; of those, \(80\%\) \texttt{[MASK]}, \(10\%\) random, \(10\%\) unchanged.
\item \textbf{GPT} = decoder-only autoregressive model using masked self-attention.
\item Encoder-only models suit representation tasks; decoder-only models suit generation; encoder-decoder models suit conditional generation such as translation.
\item \(W_Q\), \(W_K\), and \(W_V\) are shared parameters across all tokens and examples.
\item Raw attention scores come from queries against keys, not keys against values.
\item Cross-attention is specifically the decoder weighting encoder-side representations.
\item Removing the causal mask changes the training objective; it is not harmless just because inference is left-to-right.
\item Constraining \(W_Q=W_K\) may reduce expressivity, but it does not change the \(O(n^2)\) attention cost.
\item Learned absolute positional embeddings do not naturally extrapolate to longer sequence lengths than those seen in training.
\item Types of learning: supervised, unsupervised / self-supervised, semi-supervised, multi-task, transfer learning.
\item Standard attention is \(O(n^2)\) in sequence length.
\end{itemize}

\section*{Lecture 5: Evaluation}
\begin{itemize}
\item \textbf{Intrinsic evaluation}: fit to the training-style objective, e.g. held-out log-likelihood or perplexity.
\item \textbf{Extrinsic evaluation}: performance on an actual downstream task.
\item Core warning: perplexity is useful but not sufficient.
\end{itemize}

\textbf{Held-out perplexity:}
\[
\PPL
=
\exp\left\{-\frac{1}{n}\sum_{(x,y)} \log p_\theta(y \mid x)\right\}.
\]
\begin{symdefs}
\sym{\(n\)}{number of evaluation examples or tokens, depending on setup}
\end{symdefs}

\textbf{LM3 counterexample:}
\[
\begin{aligned}
p_3(x)&=\frac{1}{2}p_1(x) \\
\Longrightarrow \quad
\frac{1}{n}\sum_x \log p_3(x)
&=
\frac{1}{n}\sum_x \log p_1(x)+\log\frac{1}{2}.
\end{aligned}
\]
A constant log-likelihood penalty can hide catastrophic generation failures.

\textbf{Forced-choice accuracy:}
\[
\operatorname{Accuracy}
=
\frac{\#(\text{correct})}{\#(\text{examples})}.
\]

\begin{itemize}
\item Demo \(\neq\) solved problem.
\item Recognition and generation are different skills.
\item Forced-choice evaluations are clean and easy to score, but they test a narrow recognition setup rather than open-ended generation.
\item Free-response evaluation better matches real use, but it is harder to score and easier to make noisy.
\item Human evaluation is expensive, noisy, and hard to reproduce, but still the gold standard in many settings.
\item Elicitation matters: prompt wording, sampling choices, and rubric design can change the measured result.
\item \textbf{BLEU}: precision-oriented \(n\)-gram overlap with a brevity penalty.
\item \textbf{ROUGE}: recall-oriented overlap, common in summarization.
\item \textbf{BERTScore}: embedding-based semantic similarity rather than exact overlap.
\item LM-as-a-judge can work because verification is often easier than generation.
\item Common failure modes: Goodhart's Law, Clever Hans effects, leakage, spurious shortcuts, metric overfitting.
\end{itemize}

\section*{Lecture 6: Pretraining and Fine-Tuning}
\textbf{Transfer-learning pipeline:}
\[
\theta_{\mathrm{PT}}
=
\arg\max_\theta \sum_{(x,y)\in D_{\mathrm{PT}}}\log p_\theta(y \mid x),
\]
\[
\theta_{\mathrm{FT}}
=
\arg\max_\theta \sum_{(x,y)\in D_{\mathrm{FT}}}\log p_\theta(y \mid x)
\quad \text{initialized at } \theta_{\mathrm{PT}}.
\]
\begin{symdefs}
\sym{\(D_{\mathrm{PT}}\)}{pretraining corpus}
\sym{\(D_{\mathrm{FT}}\)}{downstream fine-tuning dataset}
\end{symdefs}

\textbf{Gradient-view intuition:}
\[
\theta' \leftarrow \theta + \nabla_\theta \log p(y_{\mathrm{PT}} \mid x_{\mathrm{PT}}),
\]
\[
\theta'' \leftarrow \theta' + \nabla_{\theta'} \log p(y_{\mathrm{FT}} \mid x_{\mathrm{FT}}).
\]

\textbf{Pretraining objectives:}
\begin{itemize}
\item autoregressive next-token prediction:
\[
p_\theta(x_{1:T})
=
\prod_{t=1}^{T} p_\theta(x_t \mid x_{1:t-1});
\]
\item masked language modeling:
\[
\max_\theta \sum_{t \in \mathcal{M}} \log p_\theta(x_t \mid x_{\setminus \mathcal{M}});
\]
\item infilling:
\[
\max_\theta \log p_\theta(x_{\mathrm{missing}} \mid x_{\mathrm{prefix}}, x_{\mathrm{suffix}});
\]
\item ranking / discrimination:
\[
p_\theta(y=1 \mid s_1, s_2).
\]
\end{itemize}

\begin{itemize}
\item Pretraining learns general language/world knowledge; fine-tuning steers behavior to tasks or domains.
\item Instruction tuning makes a base model respond usefully to prompts.
\item Continued pretraining on in-domain unlabeled text is often helpful.
\item Data mixing or reweighting tries to make pretraining data look more like the downstream domain.
\item Synthetic question-answer (QA) or instruction data can help when labels are scarce.
\item Parameter-Efficient Fine-Tuning (PEFT) updates only a small subset of parameters.
\item Fine-tuning is often \textbf{shallow}: it surfaces or reweights capabilities already present from pretraining.
\item Refusal is not the same as forgetting; unlearning is fragile.
\item Overaggressive fact fine-tuning can increase hallucination.
\end{itemize}

\section*{Lecture 7: Scaling Laws and Emergent Behavior}
\textbf{Validation loss and perplexity:}
\[
\PPL=e^{L}.
\]

\textbf{Generic scaling law:}
\[
L(X)\propto X^{-\alpha},
\qquad 0<\alpha<1.
\]

\textbf{Chinchilla-style model:}
\[
L(N,D)=E+\frac{A}{N^\alpha}+\frac{B}{D^\beta}.
\]
\begin{symdefs}
\sym{\(N\)}{parameter count}
\sym{\(D\)}{training-token count}
\sym{\(E\)}{irreducible loss floor}
\sym{\(A,B,\alpha,\beta\)}{fitted scaling-law constants}
\end{symdefs}

\textbf{Dense-Transformer compute approximation:}
\[
C \approx 6ND.
\]
\begin{symdefs}
\sym{\(C\)}{rough dense-Transformer training-compute budget}
\end{symdefs}

\textbf{Under fixed compute:}
\[
D=\frac{C}{6N}
\]
and choose \(N,D\) to minimize the predicted loss.

\textbf{Log-log linearization:}
\[
L(X)=kX^{-\alpha}
\quad \Longrightarrow \quad
\log L(X)=\log k - \alpha \log X.
\]

\begin{itemize}
\item Workflow: train many smaller models, measure loss, fit power laws, extrapolate.
\item On log-log plots, approximate power laws look roughly linear.
\item Under fixed compute, training cost is roughly proportional to parameter count times training-token count, so model size and data size trade off against each other.
\item Kaplan view: bigger models looked especially powerful.
\item Chinchilla refinement: many large models were undertrained, so compute-optimal scaling uses a better parameter/data balance.
\item In-context learning: adapt from instructions/examples in the prompt, without gradient updates.
\item Chain-of-thought reasoning: generate intermediate reasoning steps token by token.
\item Emergence may be real, but some ``jumps'' are metric artifacts.
\item Loss and perplexity remain useful but imperfect proxies for downstream reasoning and task performance.
\end{itemize}

\section*{Likely Quiz Habits}
\begin{itemize}
\item Be able to write the core equations from memory.
\item Be ready to explain \emph{why} a model/objective does or does not define a valid language model.
\item Expect compare/contrast questions: CBOW vs RNN, RNN vs Transformer, BERT vs GPT, intrinsic vs extrinsic evaluation, Kaplan vs Chinchilla.
\item When asked about limitations, mention order sensitivity, context length, optimization, parallelism, metric mismatch, and data/compute tradeoffs.
\end{itemize}

\end{multicols*}

\end{document}
